\chapter{Conclusions and Prospects}
\label{chp:conclusions-and-prospects}

\section{Conclusions}

In chapter 2, we addressed the critical gap of incomplete miRNA transcription start site annotation by developing miRStart 2.0, a deep learning platform integrating over 4,500 multi-modal sequencing datasets. The model achieved 83-93\% accuracy in predicting experimentally validated TSSs, substantially outperforming existing methods. This enabled systematic annotation of 28,828 putative TSSs across human and mouse miRNomes and mapping of over 6 million tissue-specific TF-miRNA interactions. By establishing the upstream regulatory layer of miRNA networks, this work provides the essential data foundation that enables comprehensive network analyses, directly supporting the integrative approaches pursued in subsequent modules.

Chapter 3 established the value of edge-centric network analysis through two complementary case studies. In the first study, we demonstrated that MTI-based prognostic modeling significantly outperforms expression-based approaches in triple-negative breast cancer, with superior accuracy and robustness across independent validation cohorts. This validates the clinical utility of individual edge information. In the second study, we systematically characterized hub miR-483 across six disease categories, revealing that context-specific network topology --- shared backbone with disease-specific dominant wiring --- explains its paradoxical functional duality. Together, these studies establish an analytical paradigm where edge information provides unique insights that node-centric approaches cannot capture.

Chapter 4 extended the edge-centric paradigm to the dynamic dimension of epitranscriptomic modification. We developed a computational pipeline enabling detection of 8-oxo-guanine modifications from standard miRNA-seq data, validated against gold-standard immunoprecipitation data with over 80\% concordance. Applying this pipeline to liver cancer, we introduced the Rewiring Activity Score to quantify modification-driven target switching. We identified stage-specific prognostic biomarkers and uncovered mechanistic links between oxidative rewiring and OXPHOS metabolic reprogramming. This work reveals that miRNA regulatory networks possess dynamic plasticity, with chemical modifications providing a mechanism for rapid network reconfiguration without changes in expression level.

\section{Future Directions}

The research presented in this thesis constructs a multi-dimensional framework for understanding miRNA regulation, traversing from transcriptional origins to dynamic epitranscriptomic modifications. Building upon this foundation, future investigations will expand into three critical dimensions: resolution, scale, and therapeutic translation.

First, to enhance the granularity of upstream regulation, future iterations of~miRStart~will integrate emerging~single-cell ATAC-seq (scATAC-seq)~data. Moving beyond bulk-tissue averages, this approach will decipher cell-type-specific TF-miRNA regulatory circuits, unraveling how miRNA biogenesis is differentially controlled within the complex heterogeneity of the tumor microenvironment versus healthy tissues.

Second, to scale the systematic discovery of dynamic rewiring, I aim to establish a fully automated, high-throughput computational pipeline to mine raw small-RNA sequencing data from public repositories (e.g., TCGA, GEO). This initiative will culminate in the construction of a~Global Atlas of miRNA Oxidative Modifications. By systematizing the prediction of "Oxi-miR-target" interactions, this resource will enable the comprehensive assessment of how oxidative rewiring reshapes the transcriptome across diverse human pathologies.

Finally, the translational potential of oxidative rewiring will be a primary focus. Building on the observed metabolic reprogramming (e.g., the OXPHOS shift), future studies will explore the emerging field of~"Pharmaco-Epitranscriptomics."~Specifically, we will investigate whether antioxidant interventions or specific small molecules can effectively reverse miRNA oxidation in vivo. This research aims to identify therapeutic agents capable of "resetting" the rewired regulatory networks, offering a novel strategy to combat metabolic plasticity and therapy resistance in aggressive cancers.
